\PassOptionsToPackage{brazilian}{babel}
\documentclass[
	12pt,
	a4paper,
	article,
	twocolumn
]{abntex2}

% --- Pacotes Básicos e Hifenação em Português ---
\usepackage[brazilian]{babel}
\usepackage[utf8]{inputenc}
\usepackage[T1]{fontenc}
\usepackage{lmodern}
\usepackage{indentfirst}
\usepackage{graphicx}
\graphicspath{{figuras/}{./}}
\usepackage{amsmath,amssymb}
\usepackage{booktabs}
\usepackage{enumitem}
\usepackage{hyperref}

\hypersetup{
	colorlinks=true,
	linkcolor=blue,
	citecolor=blue,
	urlcolor=blue
}

% Remove o título em língua estrangeira padrão do abntex2 article
\renewcommand{\maketitlehookb}{}

\titulo{Mini-projeto 3: Atenção e Transformers em Séries Temporais Climáticas}

\autor{
	Eduardo N. S. R.\\[0.2em]
	\small PPGCO / PGC112 -- Faculdade de Computação (FACOM)\\
	\small Universidade Federal de Uberlândia (UFU)\\
	\small Repositório GitHub: \href{https://github.com/vndmtrx/ppgco-pgc112-ia}{\texttt{vndmtrx/ppgco-pgc112-ia}}
}

\data{Uberlândia -- MG, Brasil, 2026}

\begin{document}
\frenchspacing
\SingleSpacing

\maketitle

\begin{abstract}
Este resumo será redigido após a conclusão dos experimentos, sintetizando os principais achados, métricas de previsão e conclusões sobre o mecanismo de atenção em dados climáticos.
\end{abstract}

\section{Introdução e Dados}

Esta seção será preenchida com a contextualização do conjunto de dados Jena Climate (2009--2016), a análise estatística exploratória das 14 variáveis meteorológicas, a justificativa da divisão cronológica entre treino e teste e a estrutura das janelas deslizantes de tempo.

\section{Representação Densa e Arquitetura Transformer}

Esta seção será preenchida com o detalhamento dos componentes do modelo: a projeção linear dos tokens multivariados, a formulação da codificação posicional, a arquitetura dos blocos codificadores de autoatenção (\textit{Multi-Head Attention}) e a definição fundamentada da função de perda.

\section{Previsão Multivariada e Resultados}

Esta seção será preenchida com as curvas de perda por época e a avaliação experimental das predições para os dois objetivos propostos: previsão do próximo passo ($t+1$) e previsão contínua do horizonte estendido de 24 horas.

\section{Mapas de Atenção e Interpretabilidade}

Esta seção será preenchida com a extração dos pesos de atenção das cabeças do codificador, a apresentação dos mapas de calor (\textit{heatmaps}) médios e a discussão física e meteorológica dos padrões temporais identificados.

\section{Conclusão}

Esta seção será preenchida com as considerações finais, limitações observadas e conclusões do trabalho.

\begin{thebibliography}{99}

\bibitem{vaswani2017}
VASWANI, A. et al. Attention is all you need. In: \textbf{Advances in Neural Information Processing Systems (NeurIPS)}, v.~30, 2017.

\bibitem{smola2023}
ZHANG, A.; LIPTON, Z.~C.; LI, M.; SMOLA, A.~J. \textbf{Dive into Deep Learning}. Cambridge University Press, 2023. Disponível em: \url{https://d2l.ai/}.

\end{thebibliography}

\end{document}
